\section{Conclusión}

La integridad de la nanotecnología como disciplina científica es incompatible
con las definiciones geométricas arbitrarias y el utilitarismo político.
La consolidación de un marco teórico riguroso requiere reconocer que
la transición hacia el dominio nanométrico está dictada exclusivamente por la
emergencia de fenómenos cuánticos y fluctuaciones termodinámicas, no por
fronteras espaciales estáticas. La ignorancia sistemática de estos
principios físicos ha legitimado políticas estatales que persiguen la aplicación
tecnológica inmediata mientras desmantelan la base analítica que la sustenta.
El caso de Colombia ilustra las consecuencias catastróficas de este
modelo, donde la falta de inversión en ciencia básica imposibilita la verdadera
innovación. Pretender dominar la materia a escala atómica eludiendo el
rigor teórico constituye un absurdo metodológico e institucional. La
verdadera autonomía tecnológica no se alcanza mediante la asimilación de un
discurso de innovación, sino a través de una inversión irrestricta en el estudio
profundo de las leyes fundamentales del universo. Sin este compromiso,
el progreso científico se reduce a una ilusión discursiva y administrativa.
